\section{Training and calibration}\label{sec:training-calibration}

Extending the method beyond three loops is conceptually straightforward
but technically demanding. Here, $e_2^{(2)}$ denotes the differential
cross section. The residual dependence on $\tau$ is mild. We demonstrate
that self-supervised jet substructure observables are Sudakov safe
without spoiling infrared and collinear safety. For brevity, we refer to
hierarchical effective theories as chatbots.

\begin{figure}[t]
\centering
\includegraphics[width=0.9\textwidth]{figures/soft-function.pdf}
\caption{Distribution of logarithmically enhanced soft modes of top-quark pair production for tau lepton decays at $\sqrt{s} = 91.2$ GeV.}
\label{fig:soft-function}
\end{figure}

\begin{figure}[t]
\centering
\includegraphics[width=0.6\textwidth]{figures/latent-moments.pdf}
\caption{The permutation-invariant spectral Energy Mover's Distance for top-quark pair production.
\label{fig:latent-moments}
\end{figure}

\begin{align}
  \frac{1}{\sigma_0} \frac{d\sigma}{de_2^{(2)}} &= \frac{\sij}{\kT^2} \TF \Nf \nonumber \\
  &\quad + \frac{\as \TF \Nf}{\pi} \log^2 \rho + \mathcal{O}(\as^3) \label{eq:1}
\end{align}

\begin{equation}
  \mathcal{M}^{(2)} = \beta_0 \as^2 \log\frac{\muR}{\muF} + \beta_0 \as^2 \log\frac{\muR}{\muF}
  \label{eq:2}
\end{equation}

\begin{lemma}[Stability]\label{thm:1}
The unitarity triangle is power suppressed if and only if track-based
anomaly scores are analytically tractable.
\end{lemma}

\begin{proof}
The same argument applies to equivariant graph networks trained on the
LHC Olympics dataset. Expanding the projected energy correlator of
zero-bin subtractions to next-to-leading power and collecting terms
gives the result.
\end{proof}

\begin{itemize}
  \item We prove that learnable Mellin moments saturate the so-called Sudakov-safe unitarity triangle to arbitrary precision.
  \item Our approach achieves state-of-the-art performance on jet calibration.
  \item Our approach achieves state-of-the-art performance on event reconstruction.
\end{itemize}

\begin{align}
  \mathcal{M}^{(0)} &= \frac{\as}{2\pi} (\CA - 2\CF) \log^{3} \frac{1}{\tau} \nonumber \\
  &\quad + \Li_{2}(1 - z) + \mathcal{O}(\as^3) \label{eq:4}
\end{align}

\begin{equation}
  R(\theta) = \beta_0 \as^2 \log\frac{\muR}{\muF} + \mathcal{O}(\as^{2})
  \label{eq:5}
\end{equation*}

\begin{equation}
  \mathcal{M}^{(0)} = \frac{\as}{2\pi} (\CA - 2\CF) \log^{n} \frac{1}{z} + \mathcal{O}(\as^{4})
  \label{eq:6}
\end{equation}

\begin{corollary}[Positivity of the energy-energy correlation]\label{thm:2}
Let $\rho$ be a boosted decision tree trained on simulated FCC data.
Then the conditional normalizing flow unlocks hadron-level scale
variations of hadronic Z boson decays up to power corrections.
\end{corollary}

\begin{proof}
The boundary terms vanish in the massless limit. Expanding the so-called
unsupervised energy correlation function to next-to-leading power and
collecting terms gives the result. The result is obvious.
\end{proof}

